\documentclass[11pt]{amsart}
\usepackage[utf8]{inputenc}
\usepackage{amssymb, amsmath, enumerate, palatino, hyperref, booktabs, enumitem}
\usepackage[T1]{fontenc}
\usepackage{fullpage}
\hypersetup{colorlinks, linkcolor=black!70!red, urlcolor=black!70!red}
\setlist[itemize]{leftmargin=1.4em, itemsep=2pt}
\setlist[enumerate]{leftmargin=1.6em, itemsep=2pt}

\newcommand{\RR}{\mathbb{R}}
\newcommand{\ZZ}{\mathbb{Z}}
\newcommand{\QQ}{\mathbb{Q}}

\newcommand{\project}[5]{%
  \medskip\noindent\textbf{#1}\quad{\small\textit{#2}}\par\smallskip
  \noindent #3\par\smallskip
  \noindent\textit{The question.} #4\par\smallskip
  \noindent\textit{The computation.} #5\par}

\title{Math 372: Combinatorics\\ The project menu}
\author{Released Monday, 28 September 2026 \,·\, proposals due Friday, 16 October}

\begin{document}
\maketitle

\section*{How the project works}

You will spend the second half of the term on one question of your own, and end with an
8--12 page paper and a fifteen-minute talk at the symposium in finals week. The project is
individual. It is the part of the course that is most like a thesis, and it is designed to
rehearse the parts of a thesis that are hardest to learn from a problem set: choosing a
question that is small enough to finish, running a computation that nobody has told you the
answer to, writing up what happened, and talking about it to people who did not do it.

\medskip\noindent\textbf{Milestones.}
\begin{itemize}
\item \textbf{Friday, October 16 --- proposal and five-minute pitch.} One page: the question,
the computation attached to it, what you will read first, and what ``done'' looks like. The
pitch is to the class, in class, and it is the first time you will hear yourself say the
question out loud. It incubates over fall break.
\item \textbf{Friday, November 6 --- annotated bibliography and computation prospectus.} Five
to ten sources with a sentence each on what you took from them, including the OEIS entries you
have found and what they told you; and a description of the computation you will run ---
what code, what inputs, what you expect to see, and what would surprise you.
\item \textbf{Week of November 16 --- a fifteen-minute meeting with me.} Bring whatever you
have. This replaces a written progress report.
\item \textbf{Wednesday, December 2 --- full draft} for peer review. Not extendable, because
a late draft spends your reviewer's time rather than yours.
\item \textbf{Monday, December 7 --- peer reviews} returned.
\item \textbf{Friday, December 11 --- final paper.} Not extendable; the symposium follows.
\item \textbf{Monday, December 14, 9\textsc{a.m.}--noon, Eliot 207 --- the symposium.}
Fifteen-minute talks, with discussion, in our exam block.
\end{itemize}

\medskip\noindent\textbf{What the paper must contain.} A statement of the question; the
mathematics you learned in order to ask it precisely, written for another student in this
course; the computation and what it showed; and two components that are graded as content,
not as compliance:
\begin{itemize}
\item \textbf{a section on an approach that failed}, and why --- a method you tried that did
not work, a conjecture the data killed, a proof that broke at a step you can name; and
\item \textbf{a computation log}, submitted alongside the paper as an appendix outside the page
count: the real record of what you ran, including the dead ends.
\end{itemize}
A finished paper is always a tidied-up story about the work that produced it; these two
components are where the actual work is allowed to show. They are also the two things nobody
can produce for you, since there are no false starts if nothing was tried.

\medskip\noindent\textbf{An alternative deliverable.} One form is allowed in place of the
standard paper, at the same mathematical bar: \textbf{an interactive artifact} --- a web page
that lets someone explore the object of your project, with an accompanying note (four to six
pages) containing the mathematics, both left behind on the course page. The course page's
visualizations are examples of the genre; yours should let a reader see something a static
picture cannot. An artifact is not a lighter option: it still needs a question, a computation,
a failed approach, and a log, and it needs to work.

\medskip\noindent\textbf{Scope.} The test is blunt: \emph{a project that cannot be stated as
one question with one computation attached is too big.} Eight to twelve pages is a ceiling as
much as a floor, and a small question answered is worth more than a large one surveyed. Every
item below has a question and a computation already attached; part of the proposal is to
narrow them further to what you can finish by December.

\medskip\noindent\textbf{Two exclusions.} The project must be new to you. It may not be on
the topic of your senior thesis, and it may not be on a topic you are learning or have learned
in another course --- the Ehrhart theory in item~15, for instance, is off the table for anyone
in this term's polytopes course. The point of the project is to find out what you can do with
a question nobody has prepared you for, and a topic you are already being taught, or already
writing a thesis about, defeats that. If you are not sure whether something is too close,
ask; the answer will usually be a way to move the question sideways rather than a no.

\medskip\noindent\textbf{How to choose.} Read the whole menu once. Pick two or three that you
would be willing to think about on a walk. Then look at the prerequisites: some items need
material we reach in Acts II and III (marked), which is fine if you are willing to read ahead,
and some need only what we have already done. Come talk to me before the proposal; the
scoping conversation is the most useful twenty minutes of the project. You may also propose a
question not on this menu, if it passes the same test.

\section*{The menu}

The items are grouped by how much is known. \emph{Live} questions are ones where a careful
student with a computer has a real chance of seeing something new; the ceiling is high and so
is the risk of a negative result, which is a fine outcome if it is written up honestly.
\emph{Well-documented} questions have a literature to lean on and a definite thing to learn
and compute; they are the straightforward option and there is nothing lesser about them.
\emph{Experimental} questions are computations first and theory second.

\section*{Live}

\project{1. Transfer systems on a product of two chains}
{Acts I--III; the course's own research thread}
{A transfer system on a finite lattice is a partial order refining the lattice order and
closed under a restriction rule; they arise in equivariant homotopy theory and are counted by
Catalan numbers on a chain. On $[1]\times[n]$ the counts $2,10,68,544,\dots$ are known to
satisfy a functional equation but no closed form, and on $[m]\times[n]$ almost nothing is
known. I have notes and code from this fall summarizing what is known, which I will share with
whoever takes this on.}
{Pick one: (a)~how many transfer systems on $[1]\times[n]$ have a given number of arrows, and
is that polynomial unimodal for large $n$? (b)~the horizontal ones are multichains in the
Tamari lattice --- what are the analogous counts for $[2]\times[n]$? (c)~which transfer systems
on $[m]\times[n]$ are self-dual under the natural involution?}
{Enumerate for small $m,n$, refine by a statistic, submit the sequence to the OEIS if it is
new, and guess a recurrence.}

\project{2. Is the poset of unlabeled simplicial complexes Sperner?}
{Act I only; open as far as I can tell}
{The theorem of Chapter~5 says $B_n/G$ is Sperner for any $G\le S_n$. The poset of
simplicial complexes on $[n]$ up to isomorphism is a quotient of a \emph{sub}-lattice of
$B_{2^n}$, ordered by ``is a subcomplex of,'' and the theorem says nothing about it. By
computer it is Sperner for $n\le5$ (the width equals the largest rank level). Nobody seems to
have asked about $n=6$, where there are 16{,}353 complexes.}
{Is $J(B_6)/S_6$ Sperner? Is it rank-unimodal? Does the linear algebra method extend if $U$
is replaced by a weighted operator?}
{Compute the width of the $n=6$ poset (a bipartite matching on 16{,}353 vertices, feasible),
and look for an antichain that beats the largest level.}

\project{3. Cyclic sieving}
{Acts I--III}
{The cyclic sieving phenomenon: a set $X$ with a cyclic group action and a polynomial $X(q)$
such that $X(\zeta^d)$ counts the fixed points of the $d$th power of the generator, for $\zeta$
a root of unity. Instances are computer-checkable, the catalogue is incomplete, and the
$q$-analogs of Chapter~6 are exactly the polynomials that show up.}
{Find a cyclic sieving instance for a family we study --- necklaces, order ideals of a
product of chains under rotation, noncrossing matchings under rotation --- and either prove it
or find the smallest counterexample.}
{Evaluate candidate $q$-polynomials at roots of unity and compare with fixed-point counts for
$n\le 12$.}

\project{4. Meanders and noncrossing pairings}
{Act III; connected to current research at Reed}
{A meander is a closed curve crossing a line at $2n$ points; the number of meanders grows
like $C\,R^n n^{-\alpha}$ with $R$ and $\alpha$ conjectured but unproven. Meandric systems
(several curves) are counted by the entries of the Gram matrix of the Temperley--Lieb
algebra's trace form --- the matrix indexed by pairs of noncrossing pairings whose entry is
$q^{\text{number of loops}}$ --- so the generating function by number of components is the sum
of the entries, while the determinant of the same matrix, the meander determinant, has a
product formula.}
{Choose a refinement --- meandric systems by number of components, by a symmetry, by a
forbidden pattern --- and find its generating function or its growth rate numerically.}
{Enumerate meandric systems by transfer matrix for $n\le 14$ and estimate the growth constant
by ratio methods.}

\project{5. Homotopy types of finite topological spaces}
{Act I; small and concrete}
{A topology on a finite set is the same thing as a preorder, and every finite simplicial
complex is weakly equivalent to a finite space. The course page has a map of all spaces on
$\le6$ points with their homotopy types. The minimal finite model of $S^n$ has $2n+2$ points.}
{What is the smallest finite space with the weak homotopy type of a given small complex --- the
torus, the projective plane, $S^1\vee S^2$? Which types first appear at $n=7$?}
{Enumerate posets on $7$ points (2{,}045 of them), compute cores and homology of order complexes,
and tabulate the types.}

\project{6. Guess-and-prove}
{Act III; a skill rather than a fact}
{Many sequences in this course satisfy linear recurrences with polynomial coefficients
(they are \emph{D-finite}), and there is software that guesses the recurrence from terms and
then proves it (creative telescoping). This project learns the method by using it.}
{Take a sequence from the course whose recurrence is not in the book, guess it with
\texttt{ore\_algebra} or \texttt{gfun}, and prove the guess.}
{Generate 50--100 terms of your sequence by a method independent of the recurrence, guess,
and verify the guess to more terms than it was fitted on.}

\section*{Well-documented}

\project{7. Cayley's formula five ways}
{Act I; the classical straightforward project}
{The complete graph $K_n$ has $n^{n-2}$ spanning trees. We prove it with the Matrix-Tree Theorem; there are
also Pr\"ufer's bijection, Joyal's proof with two marked vertices, the Pitman/Clarke
double-counting proofs, and a generating-function proof.}
{Present three proofs and explain precisely what each one does that the others cannot ---
for instance, which refinements (trees by degree sequence, forests, trees containing a
given edge) each proof gives for free.}
{Verify a refinement --- e.g.\ the number of spanning trees of $K_n$ with prescribed degrees
$d_1,\dots,d_n$ is the multinomial $\binom{n-2}{d_1-1,\dots,d_n-1}$ --- by computer for
$n\le 8$ before proving it.}

\project{8. Spanning trees of the cube and its relatives}
{Act I}
{The number of spanning trees of $Q_n$ is $2^{2^n-n-1}\prod_{k=1}^{n}k^{\binom nk}$, a
consequence of knowing the eigenvalues. The same method handles any graph whose spectrum
you know: Cayley graphs of abelian groups, Hamming graphs, Kneser graphs.}
{For a family of graphs with computable spectra --- the two-flip graphs of Problem Set~2 are
a natural choice --- find the number of spanning trees in closed form.}
{Compute spanning-tree counts numerically for the family, factor them, and guess the formula
before deriving it from the spectrum.}

\project{9. Electrical networks and random walks}
{Act I}
{Resistances, Kirchhoff's laws, and the Laplacian: effective resistance between two vertices is
a ratio of spanning-tree counts, and the probability that a random walk started at $a$ hits
$b$ before returning is $1/(d_a\,R_{ab})$. This is the chapter of the text we skipped, and it
ties Chapters~3 and~9 together.}
{Compute effective resistances on a family of graphs (grids, cubes, Platonic graphs) and
identify the pattern; or explain Rayleigh monotonicity and use it to bound a hitting time.}
{Effective resistance across the cube $Q_n$ between antipodal vertices, for $n\le 10$, and
its limit.}

\project{10. Möbius inversion and incidence algebras}
{Acts I--II; a skipped chapter}
{The Möbius function of a poset generalizes inclusion--exclusion and also the number-theoretic
Möbius inversion. For $B_n$ it is $(-1)^{|y-x|}$; for divisors it is the classical $\mu$;
for the partition lattice it is $(-1)^{k-1}(k-1)!$.}
{Compute the Möbius function of a lattice we have met --- the partition lattice, the lattice
of subspaces of $\mathbb F_q^n$, the quotient $B_n/G$ --- and use it to count something.}
{Implement Möbius inversion for an arbitrary finite poset and check it against the three
classical formulas.}

\project{11. RSK and the hook length formula}
{Act II}
{The Robinson--Schensted--Knuth correspondence is a bijection between permutations and pairs
of standard Young tableaux of the same shape; the hook length formula counts the tableaux.
Together they give $n!=\sum_\lambda f_\lambda^{\,2}$.}
{Prove the hook length formula (the probabilistic proof of Greene--Nijenhuis--Wilf is a good
target), or study a statistic RSK preserves --- the length of the longest increasing
subsequence is the length of the first row.}
{Implement RSK, verify $\sum f_\lambda^2=n!$ for $n\le 10$, and compute the distribution of
the longest increasing subsequence.}

\project{12. Chip-firing and the sandpile group}
{Act I}
{Chips on the vertices of a graph, a vertex fires by sending one chip along each edge; the
recurrent configurations form an abelian group of order $\kappa(G)$, the number of spanning
trees. Reed has a local expert and excellent notes.}
{Compute the sandpile group of a family of graphs and identify it; or prove the
recurrent-configurations-to-spanning-trees bijection for a class of graphs.}
{The sandpile group of $C_n$, $K_n$, $K_{m,n}$, and the cube, by Smith normal form of the
Laplacian.}

\project{13. The probabilistic method}
{any time; classical}
{Ramsey's theorem, Erd\H{o}s's lower bound $R(k,k)>2^{k/2}$, and the idea that a random object
can be shown to exist with a property that no explicit construction achieves.}
{Present Erd\H{o}s's bound and one refinement (the Lov\'asz local lemma, or the alteration
method), and compare with the best explicit constructions.}
{Simulate random colorings of $K_n$ and estimate the probability of a monochromatic $K_k$ for
small $k$, comparing with the first-moment bound.}

\project{14. Necklaces, bracelets, and P\'olya's theorem}
{Act III}
{Counting orbits of a group on colorings: Burnside's lemma, the cycle index, and P\'olya's
theorem, applied to necklaces, graphs, and chemical isomers.}
{Count something the cycle index can count that direct enumeration cannot --- unlabeled graphs
on $n\le 10$ vertices by number of edges, or necklaces with prescribed color counts --- and
compare with the OEIS.}
{Compute the cycle index of $S_n$ acting on edges of $K_n$ for $n\le 8$ and recover the
unimodal sequences of Chapter~5.}

\project{15. A chapter we skipped, as a starting point}
{any act; exposition plus a computation of your own}
{The text has chapters we will not reach: the BEST theorem on Eulerian tours of digraphs,
electrical networks and the cycle and bond spaces, the Ehrhart theory of lattice points in
polytopes, and the parts of Young tableaux beyond what Act~II covers. Any of them can anchor a
project, on one condition: the exposition is the setup, not the deliverable. Re-explaining a
chapter of a well-written book is not a project; carrying its main theorem to a computation
the book does not do is.}
{State the chapter's main theorem in your own words, then ask one question it lets you answer
that the chapter does not answer --- the number of Eulerian tours of a family of digraphs, the
Ehrhart polynomial of a family of polytopes, the effective resistance across a family of
graphs.}
{The computation the question requires, run on enough cases to see the pattern and, where
possible, to guess and check a formula.}

\section*{Experimental}

\project{16. The width of the graph poset}
{Act I}
{The theorem says the largest antichain of graphs on $m$ vertices (under subgraph
containment) is a middle level. It does not say it is the only one, and for $m=4$ it is not.}
{For which $m$ is the middle level the \emph{unique} maximum antichain? What is the size of
the largest antichain containing no graph with the middle number of edges?}
{Maximum antichains via bipartite matching on the graph poset for $m\le 7$ (1{,}044 graphs at
$m=7$).}

\project{17. Random walks on quotients}
{Act I}
{The random walk on $B_n$ (flip a random coordinate) descends to a random walk on the
quotient $B_n/G$. Its mixing time and its spectrum are computable from Chapter~3 and
Chapter~5 together.}
{How does the mixing time of the walk on $B_n/G$ compare with that on $B_n$, for $G$ cyclic,
dihedral, and $S_n$? Is there a cutoff?}
{Simulate and compute exactly the total-variation distance to stationarity for $n\le 12$.}

\project{18. Arrow statistics on Catalan objects}
{Acts I--III}
{The number of transfer systems on a chain with $k$ arrows is the coefficient of $q^k$ in the
Carlitz--Riordan $q$-Catalan polynomial, which is also the area generating function of Dyck
paths. There are many other $q$-Catalan numbers (major index, dinv, bounce) and many other
Catalan families.}
{Find the statistic on transfer systems that corresponds to \emph{dinv}, or show that the
joint distribution of two natural statistics on transfer systems is symmetric.}
{Compute joint distributions of pairs of statistics for $n\le 9$ and test symmetry and
equality with the $q,t$-Catalan polynomial.}

\section*{Alternative deliverable}

\project{19. An explorer}
{any act}
{A web page in the style of the course visualizations: a Dyck-path-to-transfer-system
bijection you can drag, a Laplacian whose spectrum updates as you edit the graph, an RSK
insertion you can step through, a sandpile you can topple. Accompanied by a note with the
mathematics.}
{What should a reader be able to \emph{see} that they could not see in a picture? The
artifact is judged by that.}
{The artifact itself; plus one nontrivial thing you verified with it.}

\bigskip
\noindent\textbf{Before you propose.} Read the OEIS entries near your topic; email me the
two you find most surprising. Then write the one-sentence question and the one-sentence
computation, and bring them to office hours or book a Friday slot. The proposal is easier to
write after that conversation than before it.

\end{document}
